\documentclass{article}[11pt]
\usepackage[french]{babel}
\usepackage[utf8]{inputenc}
\usepackage{amsmath, amssymb, amsthm, mathtools}
\usepackage[margin=2.5cm]{geometry}
\usepackage{hyperref}
\usepackage{cleveref}
\usepackage{tikz, enumitem}
\usepackage{stmaryrd}
\usepackage{centernot}

\newtheorem{theorem}{Théorème}[section]
\newtheorem{proposition}[theorem]{Proposition}
\newtheorem{property}[theorem]{Propriété}
\newtheorem{corollary}[theorem]{Corollaire}
\newtheorem{lemma}[theorem]{Lemme}
\theoremstyle{definition}
\newtheorem{definition}[theorem]{Définition}
\newtheorem{remarque}[theorem]{Remarque}

\newcommand{\R}{\mathbb{R}}
\newcommand{\eps}{\varepsilon}
\newcommand{\norm}[1]{\lVert #1 \rVert}
\title{Notes de Cours - Intégrale de Lebesgue}
\author {Mohamed Aziz Gannar}
\begin{document}
\maketitle
\tableofcontents
\section{Ensembles mesurables}
\begin{definition}[Tribu ou $\sigma$-algebre]\label{tribu}
  Soit $E$ un ensemble quelconque. On appelle tribu sur E tout ensemble $\eps$ de parties de E vérifiant :
  \begin{itemize}
    \item{$E \subset \mathcal{E}$}
    \item Si $A \in \mathcal{E} \implies A^c \in \mathcal{E}$
    \item{Si $(A_n)_{n \in \mathbb{N}} $ une suite d'éléments de $\mathcal{E} \implies \bigcup_{n\in \mathbb{N}}A_n \in \mathcal{E}$}
  \end{itemize}
\end{definition}

\begin{remarque}
  Si $\mathcal{E}$ est une tribu sur $E$ alors: 
  \begin{itemize}
  \item $\emptyset \in \mathcal{E}$ puisque $E \in \mathcal{E}$ et $E^c = \emptyset$
  \item $ (A_n)_{n >= 0}$ est une suite d'éléments de $\mathcal{E}$ alors $\bigcap_{n>=0}A_n = (\bigcup_{n >=0}(A_n)^c)^c$ 
  \item Si $A,...,A_n$ sont des éléments de $\mathcal{E}$ alors $\cup_{k=1}^{n}A_k \in \mathcal{E}$
  \end{itemize}
\end{remarque} 

\begin{property}\label{intersection_de_tribu}
  Soit $E$ un ensemble quelconque, $(\mathcal{E}_i)_{i \in I}$ une famille de tribus sur $E$.
  L'ensemble des parties de $E$, 
  \[\bigcap_{i \in I} \mathcal{E}_i\]
  est une tribu sur $\mathcal{E}$.
\end{property}
\begin{proof} On montre chaque propriété de la définition d'une tribu.

  \begin{itemize}
    \item $\forall i \in I, E \in \mathcal{E}_i $ donc $ E \in \cap_{i \in I}\mathcal{E}_i$
    \item Soit $:warningA \in \bigcap_{i \in I}\mathcal{E}_i $ alors $\forall i \in I, A \in \mathcal{E}_i,$ donc $\forall i \in I, A^c \in \mathcal{E}_i$.
      Donc $A^c \in \bigcap_{i \in I} \mathcal{E}_i$
    \item Soit $(A_n)_{n \in \mathbb{N}}$ une suite d'éléments de $\bigcap_{i \in I} \mathcal{E}_i$.\newline
      Alors, pour tout $i \in I$, $(A_n)_{n\in\mathbb{N}}$ est une suite d'éléments de $\mathcal{E}_i$.\newline
      Donc, $\bigcup_{k \in \mathbb{N}}(A_k) \in \mathcal{E}_i $ car $\mathcal{E}_i$ est une tribu. \newline
      Donc, $\bigcup_{n \in \mathbb{N}}A_n \in \cap_{i \in I} \mathcal{E}_i$.
  \end{itemize}
  $\bigcup_{i \in I} \mathcal{E}_i$ est bien une tribu sur E.
\end{proof}
\begin{definition}[Tribu Engendrée]\label{tribu_engendree}
  Soit $E$ un ensemble quelconque et $\mathcal{C}$ une classe de parties de $E$. L'ensemble $$\bigcap_{\substack{\mathcal{E} \text{tribu} \\ \mathcal{C} \subset E}}\mathcal{E}$$ est une tribu sur $E$.\newline
  C'est la plus petite tribu sur $E$ contenant $\mathcal{C}$ et on la note $\sigma(\mathcal{C})$ et on l'appelle tribu engendrée par $\mathcal{C}$
\end{definition}
\begin{remarque}\label{def_borelien}
  Si $E$ est un espace topologique dont on note $\mathcal{O}$ la classe des ouverts et $\mathcal{F}$ la classe des fermés, on a :
  $$\mathcal{B}(E) = \sigma(\mathcal{F}) = \sigma(\mathcal{O})$$
\end{remarque}
\begin{proof}
  Prouvons que $\sigma(\mathcal{O}) \subset \sigma(\mathcal{F})$
  Il suffit de montrer que pour tout $ O \in \mathcal{O} \text{ on a } O \in \sigma(\mathcal{F})$
  Car alors, $\sigma(\mathcal{F})$ est une tribu sur E et contient tout les ouverts de E. Donc elle contient la plus petite tribu content $\mathcal{O}$.\newline
  Soit $O \in \mathcal{O}$, on a $O = (O^c)^c$. \newline
  $O$ étant un ouvert de E, $O^c$ est un fermé de E. Donc $(O^c)^c$ est un ouvert de E.\newline
  On a donc $O^c \in \sigma(\mathcal{F})$ donc par stabilité du passage au complémentaire, $ (O^c)^c \in \sigma(\mathcal{F})$.\newline
  Ainsi, on a bien $\sigma(\mathcal{O}) \subset \sigma(\mathcal{F})$.
  L'inclusion réciproque se montre de manière analogue.
  
\end{proof}
\section{Espaces Mesurés}
\begin{definition}[Mesure]\label{def_mesure}
  Soit $(E,\mathcal{E})$ un espace mesurable. Une application $\mu$ définie sur $\mathcal{E}$ à valeurs dans $[0,+\infty]$ est appelée mesure lorsque: 
 \begin{itemize}
  \item $\mu(\emptyset) = 0$.
  \item L'application $\mu$ est $\sigma$-additive, c'est à dire que , si $A_1,A_2,...$ est une famille dénombrable de parties de $E$ appartenant à $\mathcal{E}$ 2 à 2 disjointes, alors :
    $$\mu(\bigcup_{k=1}^{\infty}A_k) = \sum_{k=1}^{\infty}\mu(A_k)$$
 \end{itemize} 
\end{definition}
\begin{definition}[Espace Mesurés]\label{def_espace_mesuré}
  Si $E$ est un ensemble, $\mathcal{E}$ est une tribu sur $E$ et $\mu$ est une mesure sur $(E,\mathcal{E})$, on dit que $(E,\mathcal{E},\mu)$ est une espace mesuré.
\end{definition}
\begin{proposition}
  Il existe une unique mesure sure $(\mathbb{R}, \mathcal{B}(\mathbb(R)))$ traditionellement notée $\lambda$, tel que pour tout $a,b \in \mathbb{R}$
  $$\lambda(]a,b[) = b -a $$
  Cette mesure est la mesure de Lebesgue.
\end{proposition}
\begin{definition}[Types de Mesures]
  Soit $(E,\mathcal{E},\mu)$ un espace mesuré.
  \begin{itemize}
    \item Si $\mu(E) < \infty$. On dit que $\mu$ est une mesure finie.
    \item Si $\mu(E) = 1$. On dit que $\mu$ est une mesure de probabilité.
    \item Si il existe une suite $(A_n)_{n \in \mathbb{N}}$ d'éléments de $\mathcal{E}$ tel que :
      \begin{enumerate}
        \item $\displaystyle\bigcup_{n \in \mathbb{N}}(A_n) = E$
        \item $\forall n \geq 1\text{, } \mu(A_n) < \infty$
      \end{enumerate}
      On dit que $\mu$ est une mesure $\sigma$-finie.
    \item On dit que $x \in E $ est un atome de $\mu$ si :
      $$ \{x\} \in \mathcal{E} \text{ et } \mu(\{x\}) > 0$$
    \item On dit que $\mu$ est diffuse quand $\mu$ n'a pas d'atome.
  \end{itemize}
\end{definition}
\begin{property}
  Soit $(E,\mathcal{E},\mu)$ un espace mesuré.
  \begin{enumerate}
    \item Pour tous $A,B \in \mathcal{E}$, tel que $A \cap B = \emptyset$
      $$\mu (A \cup B) = \mu(A) + \mu(B)$$
    \item Pour tous $A,B \in \mathcal{E}$, on a:
      $$\mu(A) + \mu(B) = \mu(A\cup B) + \mu (A \cap B )\centernot \implies \mu(A\cup B) = \mu(A) + \mu(B) - \mu(A\cap B)$$
    \item Pour tous $A,B \in \mathcal{E} \text{ tel que } A \subset B$, on a:
      $$\mu(A) \leq \mu(B)$$
  \end{enumerate}
\end{property}
\begin{proof}
  On va démontrer les propriétés une à une:
  \begin{enumerate}
    \item Soient $A,B \in \mathcal{E} \text{ tel que } A \cap B = \emptyset$. La suite $(A_n)_{n \in \mathbb{N}}$ définie par: 
      $$ A_1 = A, A_2 = B, A_n = \emptyset$$
      est une suite d'éléments de $\mathcal{E}$ 2 à 2 disjointe. On a donc:
      $$\mu(A \cup B) = \mu(\bigcup_{n \geq 1}A_n) = \sum_{ n \geq 1}\mu(A_n) = \mu(A) + \mu(B)$$
    \item Soient $A,B \in \mathcal{E}$, on a:
      \[ A = (A \cap B) \cup (A \cap B^c) \tag{réunion de partie disjointe de E}\]
      \[ B = (B \cap A) \cup (B \cap A^c) \tag{réunion de partie disjointe de E}\]
      \[ A \cup B = (A \cap B^c )\cup(A \cap B)\cup (A^c \cap B)\]
      Donc,
      \[ \mu(A) = \mu(A\cap B) + \mu(A \cap B^c )\]
      \[ \mu(B) = \mu(B\cap A) + \mu(B \cap A^c )\]
      \[ \mu(A \cup B) = \mu(A\cap B^c) + \mu(A \cap B ) + \mu(A^c \cap B)\]
      Donc
      \begin{equation}
        \begin{split}
          \mu(A) + \mu(B) &= 2 * \mu(A \cap B) + \mu(A \cap B^c) + \mu(A^c \cap B) \\ 
                          &= \mu(A \cup B) + \mu(A \cap B)
        \end{split}
      \end{equation}
    \item Soient $A,B \in \mathcal{E}$, tel que $ A \subset B$, on a, $B = (A \cap B) \cup (A^c \cap B)$ qui est une union disjointe.\\ 
      Donc ,
      \begin{equation}
        \begin{split}
          \mu(B) &= \mu( A \cap B) + \mu( A^c \cap B)\\ 
                 &= \mu(A) + \mu(A^c \cap B) \geq \mu(A) 
        \end{split}
      \end{equation}
      car $ \mu( A^c \cap B) \geq 0$.
  \end{enumerate}
\end{proof}
\begin{proposition}[Propriété de continuité croissante des mesures positives]
  Soient $(E, \mathcal{E}, \mu)$ un espace mesuré et $(A_n)_{n \geq 0}$ une suite croissante d'éléments de $\mathcal{E}$.
  \[ \lim_{n\to\infty} \mu(A_n) = \sup_{n \geq 0} \mu(A_n) = \mu(\bigcup_{n \geq 0} A_n)\]
\end{proposition}
\begin{proof}
  Soit $(A_n)_{n \geq 0}$ une suite d'éléments de $\mathcal{E}$.
  D'après la propriété précédente : $(\mu(A_n))_{n \geq 0}$ est une suite croissante d'éléments de $\overline{\mathbb{R}}$ donc: 
  \[ \lim_{n\to\infty}\mu(A_n) = \sup_{n \geq 0} \mu(A_n)\]
  On pose $B_0 = A_0$ et pour tout $ n\geq 1 \text{, } B_n = A_n \backslash A_{n-1} = A_n \cap (A_{n-1})^c$
  La suite $(B_n)_{n \geq 0}$ ayant les propriétés :
  \begin{enumerate}
    \item \begin{itemize}
      \item les éléments de $(B_n)_{n \geq 0}$ sont 2 à 2 disjoints.
      \item $\bigcup_{n \geq 0}B_n = \bigcup_{n \geq 0}A_n$
    \end{itemize}
  \item $\displaystyle \forall n \geq 0 \text{, } \bigcup_{l = 0}^{n} B_l = A_n \text{ et } \bigcup_{l=0}^{\infty}B_l = \bigcup_{l=0}^{\infty}A_l$
  \end{enumerate}
  On a donc,
  \begin{equation}
  \begin{split}
    \mu(\bigcup_{l=0}^{\infty} A_l) &= \mu(\bigcup_{l=0}^{\infty}B_l)\\ 
                  &= \sum_{l=0}^{\infty}\mu(B_l) \\ 
                  &= \lim_{n \to \infty} \sum_{ l=0}^{\infty}B_l\\ 
                  &= \lim_{n \to \infty} \mu(\bigcup_{l=0}^{n}B_l) \\ 
                  &= \lim_{n \to\infty}\mu(A_n)
  \end{split}
  \end{equation}
  Montrons les propriétés $(i)$ et $(ii)$ :  
  \begin{enumerate}
    \item Soient $m,n \in \mathbb{N}$  tel que $n \centernot = m$, par exemple $ m < n$:
      \[ B_m \subset A_m \text{ et } B_m \subset (A_{n-1}^c) \subset (A_m)^c \]
      Puisque $A_m \subset A_{n-1}$.
      Donc $B_m \cap B_n = \emptyset$.
    \item Pour tout $n \in \mathbb{N}, l\in\llbracket 0,n\rrbracket$, on a $ B_l \subset A_l \subset A_n$
      Donc,  $$\sum_{l=0}^{\infty} B_l \subset A_n$$
      Soit $n\in\mathbb{N}, x \in A_n$
      On pose $n(x) = min\{k \in \mathbb{N}, x \in A_k\}$ et on a $x\in B_{n(x)}$ et $n(x) \leq n$ \text{ donc } $\displaystyle x\in \bigcup_{l=0}^{n}B_l \text{ donc } A_n \subset \bigcup_{l =0 }^{n}B_l$.\\ 
      \underline{Finalement,} puisque $\displaystyle\forall n \geq 0 A_n \subset \bigcup_{l=0}^{n}B_l \subset \bigcup_{l=0}^{\infty} B_l$ \\ 
      et $ B_n \subset A_n $
      \[ \forall n \in \mathbb{N}, B_n \subset \bigcup_{l=0}^{\infty}A_l \text{ donc } \bigcup_{l=0}^{\infty}B_l \subset \bigcup_{l=0}^{\infty}A_l  \]
  \end{enumerate}
\end{proof}
\begin{corollary}[Propriété de continuité décroissante des mesures positives]
  Soit $(E,\mathcal{E},\mu)$ un espace mesuré et $(A_n)_{n \geq 0}$ une suite décroissante d'éléments de $\mathcal{E}$
  tel que $\exists n_0 \in \mathbb{N} \text{ tel que } \mu(A_{n_0}) < \infty$.\\ 
  Alors, 
  \begin{equation}
    \begin{split}
      \lim_{n \to \infty} \mu(A_n) &= \inf_{n \geq 0 } \mu(A_n)  \\ 
                                   &= \mu(\bigcap_{n \geq 0}(A_n))
    \end{split}
  \end{equation}
\end{corollary}
\begin{remarque}[Contre Exemple]
  Soit $(\mathbb{R}, \mathcal{B}(\mathbb{R}), \lambda)$ un espace mesuré et on définit la suite décroissante $(A_n)_{n \in \mathbb{N}}$ par :
  $$ \forall n \in \mathbb{N}, A_n = ]- \infty,-n[$$
  On a donc :
  $$ \bigcap_{n \in \mathbb{N}} A_n = \emptyset$$
  et,
  $$\forall n \in \mathbb{N}, \lambda(A_n) = + \infty$$
  Or,
  $$\mu(\bigcap_{ n \in \mathbb{N}}A_n) = 0$$
  Donc on a : 
  $$ + \infty = \lim_{n \to \infty}\lambda(A_n) \centernot = \lambda(\bigcap_{n \in \mathbb{N}}A_n) = 0$$
\end{remarque}
\begin{proof}
  Soit $(A_n)_{n \geq 0}$ une suite décroissante d'éléments de E tel que $\exists n_0 \text{ tel que } \mu(A_{n_0}) < \infty$.  En posant pour tout $n \in \mathbb{N}, B_n = A_{n_0} \backslash A_n$
  On se donne une suite croissante d'éléments de $\mathcal{E}$ car $\forall n \in \mathbb{N} $
  \[ A_{n+1}  \subset A_n \implies (A_n)^c \subset (A_{n+1})^c \implies B_n = A_{n_0} \cap (A_n)^c \subset B_{n+1} \]
  
\end{proof}
\end{document}
